\section{Subjecthood as a legal category}
\label{sec:31.1}
Where a line has to be drawn, prefer one nobody administers. Law regularly draws decidable lines across continuous processes — viability, brain death, the age of criminal responsibility — and I defend them as the least-bad place for an enforceable line rather than as scientific facts. The usual worry is gaming, an error of over-inclusion. The opposite error is the one with an interest behind it: a capacity examination has an examiner, and an operator examining its own system breaks the own-case rule: the exam becomes an instrument for withholding status. Literacy tests were facially neutral and administered at a registrar's discretion, and \S\ref{sec:3.1} shows what the discretion did: a test that every white applicant passed and nearly every Black one failed measured the registrar and not the applicant \autocite{guinn1915,congress1965voting}. What the franchise kept when the tests went was the \emph{age} it had always had, and the virtue of an age is that nobody administers it.

The floor takes the civil and political articles and leaves the economic ones; what a bearer is owed draws on the second, and that is consistent: the first list is what a bearer refuses to do, the second what a polity owes it.
